\vspace{-0.2cm}
\section{SPO-Tree for Long CoT}
\vspace{-0.2cm}
\label{sec:tree-like-spo-for-long-cot-scenario}
For long Chain-of-Thought (CoT) scenarios, where the sampling cost of chain-based MC estimation in SPO-chain becomes prohibitively high, we propose a novel tree-based segment advantage estimation method. This approach drastically reduces sampling overhead, making it feasible to apply our framework effectively in long CoT settings. We refer to this instantiation of our SPO framework as \textbf{SPO-tree}. In particular, SPO-tree features the fixed token count partition, tree-based segment advantage estimation, and policy gradient with probability masks.

\noindent\textbf{(1) Fixed Token Count Segment Partition.}
\label{subsec:segment-partition-based-on-fixed-token-intervals}
In long CoT scenarios, each segment typically contains a large number of tokens. As a result, it is unlikely that all token transitions within an entire segment will have probabilities close to 1. Thus, we employ a partition strategy where segment boundaries are made at fixed token intervals, which allows us to generate segments with equal token count, supporting our tree-based segment advantage estimation method proposed in the following subsection.

\noindent\textbf{(2) Tree-based Segment Advantage Estimation.}
\label{subsec:tree-like-segment-advantage-estimation}
The core drawback of the chain-based segment advantage estimation strategy studied in Section \ref{subsec:chain-like-segment-advantage-estimaiton} is that it discards samples used for estimating the values (e.g., $\boldsymbol \tau^{(1)}$, $\boldsymbol \tau^{(2)}$, $\boldsymbol \tau^{(3)}$ in Figure \hyperref[fig:sampling_method]{2(a)}), leading to substantial waste of samples in scenarios involving long reasoning trajectories. To address this issue, we propose a tree-based segment advantage estimation strategy, which reuses the samples employed for value estimation in policy optimization, significantly improving the sample efficiency.

As illustrated in Figure~\hyperref[fig:sampling_method]{2(b)}, we model the trajectory sampling process as a tree structure. We define \(\operatorname{hist}(n)\) as the full history sequence corresponding to node \( n \), including the initial prompt \(\boldsymbol{x}\) and all tokens generated along the path from the root node to node \( n \). Each node \( n \) represents a segment \(\operatorname{seg}(n)\), which is generated by extending the sequence of its parent node, \(\operatorname{Pa}(n)\), with \(M\) newly sampled tokens, i.e., 
\[
\operatorname{seg}(n) = \big[y^{(n)}_1, \dots, y^{(n)}_{M}\big], \quad y^{(n)}_t \sim \pi\big(\cdot\,\big|\,\big[\operatorname{hist}(\operatorname{Pa}(n)), \boldsymbol y^{(n)}_{<t}\big]\big).
\]

Each node \(n\) generates a set of child nodes, denoted as \(\operatorname{Ch}(n)\), and has a set of siblings, denoted as \(\operatorname{Sib}(n)\). Nodes among siblings share the same prompts and sequence lengths (except for leaf nodes), ensuring fair comparison under the same token budget. The value of the state represented by node \( n \) is denoted as \(V(n)\), which can be estimated recursively from leaves to root as follows:
\[
\hat{V}(n) =
\begin{cases}
\mathcal{R}(\boldsymbol{x}, \operatorname{hist}(n)), & \text{if } n \text{ is a leaf node} \\
\frac{1}{|\operatorname{Ch}(n)|} \sum_{n' \in \operatorname{Ch}(n)} \hat{V}(n'), & \text{otherwise}
\end{cases}.
\]
We denote the advantage of the segment represented by node \( n \) as \( A(n) \). The estimated advantage \( \hat A(n) \) for node \( n \), relative to its siblings, is computed either in unnormalized or normalized form as:
\[ 
\hat{A}(n)= \hat{V}(n)-\texttt{mean}_{n'\in \operatorname{Sib}(n)}\hat{V}(n') \quad \text{ or } \quad  \frac{\hat{V}(n)-\texttt{mean}_{n'\in \operatorname{Sib}(n)}\hat{V}(n')}{\texttt{std}_{n'\in \operatorname{Sib}(n)}\hat{V}(n')}. 
\]
More details about the tree construction are presented in Appendix \ref{appendix:pseudo-code-of-tree-structured-advantage-estimation-method}. In contrast to the chain-based segment advantage estimation strategy, each node (segment) in the tree can be used as a training example, substantially improving sample efficiency. In addition, due to extensive node-sharing among trajectories, the actual number of tokens we need to sample is significantly less than the sum of tokens across all trajectories, greatly reducing sampling overhead. Within each training iteration, sampling a large number of trajectories from a single question can lead to the model overfitting to specific samples. To address this, we introduce a replay buffer mechanism that distributes sampled trajectories across multiple iterations in Appendix \ref{appendix:replay_buffer}.

The width of the tree determines the number of sibling nodes in each group, with a larger width providing more nodes for comparison and thus more nuanced advantage estimates. The depth of the tree controls the granularity of partitions; deeper trees yield finer-grained signals. Furthermore, the hierarchical tree structure naturally aligns with the varying uncertainty along the reasoning trajectory. Earlier segments (closer to the root), which have higher uncertainty, aggregate more samples for estimating their \(V\) values, while later segments (closer to the leaves), which are more certain, aggregate fewer samples.

\noindent\textbf{(3) Policy Gradient with Token-Probability Masks.}
\label{subsec:policy-optimization-using-edges-with-non-zero-advantage}
To focus model training on segments with discriminative signals, we extract segments with non-zero advantages from the tree for training. We also note that such filtering of non-zero advantage samples aligns with practices in several concurrent works \cite{yu2025dapoopensourcellmreinforcement, lin2025cppoacceleratingtraininggroup}. Thus, we propose to minimize \(\mathcal{J}_{\text{SPO}}^{\text{tree}}\) below for policy optimization: 
\begin{equation}
\begin{aligned}\label{eq:spo_tree}
&\mathcal J_{\text{SPO}}^{\operatorname{tree}}(\theta)= 
\mathbb E_{\substack{\boldsymbol x\sim\mathcal D,\,\text{tree}\sim\pi_{\theta_{\text{old}}}}}\Bigg\{
\frac{1}{Z}
\sum_{n\in\operatorname{tree}: \hat A(n)\neq 0}
% \sum_{\substack{n\in\operatorname{tree}\\ \hat A(n)\neq 0}}
\sum_{t=1}^{|\operatorname{seg}(n)|}\\[-0.5em]
&\qquad\quad \bigg[M_{n,t}\cdot\min\Big(r_{n,t}(\theta)\hat A(n),\,\operatorname{clip}(r_{n,t}(\theta),1-\epsilon,1+\epsilon)\hat A(n)\Big) 
-\beta D_{\mathrm{KL}}(\pi_{\theta}\|\pi_{\text{ref}})\bigg]
\Bigg\},
\end{aligned}
\end{equation}
where \(M_{n,t}\) is the token-probability mask whose value is \(1\) if \(
\pi_{\theta}(y^{(n)}_t~|~ [\operatorname{hist}(\operatorname{Pa}(n)), \boldsymbol y^{(n)}_{<t}]) < \rho\) and \(0\) otherwise, i.e.,
\(
M_{n,t}:= \mathbb{I}\Big(\pi_{\theta}(y^{(n)}_t~|~ [\operatorname{hist}(\operatorname{Pa}(n)), \boldsymbol y^{(n)}_{<t}]) < \rho\Big).
\)
The ratio term \(r_{n,t}(\theta)\) is then defined as \(\frac{\pi_{\theta}(y^{(n)}_t|[\operatorname{hist}(\operatorname{Pa}(n)),\boldsymbol y^{(n)}_{<t}])}{\pi_{\theta_{\text{old}}}(y^{(n)}_t|[\operatorname{hist}(\operatorname{Pa}(n)),\boldsymbol y^{(n)}_{<t}])}\), and the normalization term \(Z\) is given by \(Z = \sum_{n\in \text{Tree}:\,\hat A(n)\neq 0}\sum_{t=1}^{|\operatorname{seg}(n)|}M_{n,t}\).